% Generated draft: requires mathematical and source review.
\documentclass[11pt]{article}
\usepackage[margin=25mm]{geometry}
\usepackage{fontspec}
\setmainfont{Latin Modern Roman}
\usepackage{amsmath,amssymb,mathtools,unicode-math}
\setmathfont{Latin Modern Math}
\usepackage{enumitem,needspace,xurl,hyperref}
\setlength{\parindent}{0pt}
\setlength{\parskip}{5pt}
\setlength{\emergencystretch}{3em}
\newcommand{\N}{\mathbb{N}}
\newcommand{\Z}{\mathbb{Z}}
\newcommand{\Q}{\mathbb{Q}}
\newcommand{\R}{\mathbb{R}}
\newcommand{\C}{\mathbb{C}}
\newcommand{\F}{\mathbb{F}}
\newcommand{\A}{\mathbb{A}}
\newcommand{\PP}{\mathbb{P}}
\newcommand{\E}{\mathbb{E}}
\newcommand{\eps}{\varepsilon}
\newcommand{\dd}{\,\mathrm{d}}
\begin{document}
\section*{UnsolvedMath ID: 800610}
\textbf{Source problem code:} catalogue entry
\subsection{Definitions and mathematical statement}
For a topological space \(X\), let \(C_p(X)=C(X,\R)\) denote the space of continuous real-valued functions with the topology of pointwise convergence, and let \(\mathbf 0\) be the constant zero function. A basic neighborhood of \(\mathbf 0\) is
\[
N(F,\varepsilon)=\{f\in C(X,\R): |f(x)|<\varepsilon\text{ for every }x\in F\},
\]
where \(F\subseteq X\) is finite and \(\varepsilon>0\).

Under the conventional definitions used for this draft, let
\[
\Omega_{\mathbf 0}=\{A\subseteq C_p(X)\setminus\{\mathbf 0\}:\mathbf 0\in\overline A\},
\]
and let \(\Gamma_{\mathbf 0}\) consist of infinite sets \(A\subseteq C_p(X)\setminus\{\mathbf 0\}\) such that every neighborhood of \(\mathbf 0\) contains all but finitely many elements of \(A\).

An open \(\omega\)-cover of \(X\) is a family \(\mathcal U\) of open subsets of \(X\) such that \(X\notin\mathcal U\) and every finite subset of \(X\) is contained in some member of \(\mathcal U\). An open \(\gamma\)-cover is an infinite open cover \(\mathcal U\) with \(X\notin\mathcal U\) such that each point of \(X\) belongs to all but finitely many members of \(\mathcal U\). Write \(\Omega(X)\) and \(\Gamma(X)\) for these respective families of covers.

For families \(\mathcal A,\mathcal B\) of sets, the principle \(\mathsf S_1(\mathcal A,\mathcal B)\) means that for every sequence \((A_n)_{n\geq 1}\) with \(A_n\in\mathcal A\), there exist choices \(a_n\in A_n\) for all \(n\geq 1\) such that \(\{a_n:n\geq 1\}\in\mathcal B\).

\textbf{Question (Scheepers, as attributed in the supplied record).} Is
\[
\mathsf S_1\bigl(\Gamma_{\mathbf 0},\Omega_{\mathbf 0}\bigr)
\quad\Longleftrightarrow\quad
\mathsf S_1\bigl(\Gamma(X),\Omega(X)\bigr)
\]
for \(C_p(X)\) and \(X\), respectively?

\textbf{Unresolved scope.} The intended assertion would quantify over every \(X\) in the source's intended class of spaces. The supplied record does not identify that class or reproduce the source's standing conventions. In particular, it does not establish whether \(X\) must be a set of reals, a Tychonoff space, or satisfy further hypotheses. The definitions above provide an explicit conventional interpretation, but their agreement with the source, including any countability restrictions on covers or convergent sets, requires review. Thus the intended universally quantified formulation remains conditional on recovering those conventions.
\subsection{Short English statement}
The question asks whether selecting one function from each of a sequence of sets converging to the zero function, so that the selected functions accumulate at zero, is equivalent to selecting one member from each of a sequence of open \(\gamma\)-covers of \(X\) so that the selected members form an \(\omega\)-cover. The intended class of spaces and the source's standing conventions are not specified in the supplied record.
\subsection{Sources}
\begin{itemize}
\item Supplied local UnsolvedMath registry record: permanent ID \(800610\), source problem code \texttt{catalogue entry}, title ``Selection Principles in Mathematics --- Problem 38 (Scheepers).'' The registry lists Tsaban as proposer and 2006 as the proposed year; the background attributes the problem to Scheepers.
\item Tsaban, \emph{Selection Principles and Special Sets of Reals: Open Problems} (2006), Problem 38. Supplied source URL: \url{https://arxiv.org/abs/math/0609601}. The record reports access on 2026-10-16 and extraction by \texttt{arxiv\_latex}; these are supplied provenance details, not a claim of independent source access or verification.
\item The registry labels the problem \texttt{open} and records an update dated 2026-10-16. This is supplied status metadata, not an independently verified assessment of the literature.
\end{itemize}
\end{document}
